\documentclass[11pt]{article}
\usepackage[utf8]{inputenc}
\usepackage[T1]{fontenc}
\usepackage{lmodern}
\usepackage[margin=1in]{geometry}
\usepackage{booktabs}
\usepackage{parskip}
\usepackage{microtype}
\usepackage[hidelinks]{hyperref}
\setlength{\emergencystretch}{2em}
\begin{document}
\begin{center}{\LARGE\bfseries Preregistration: Heritability and selection of strategic traits in the NFL coaching tree}\end{center}
\vspace{0.5em}

Preregistration of new analyses of an existing dataset. Drafted 2026-07-01 by Jon Williamson.

\section*{0. Study type and honest-disclosure statement}

This work began as exploratory analysis of an existing dataset. As it matured into a set of confirmatory hypotheses, prompted by Mesoudi's (2020) study ``Cultural evolution of football tactics: strategic social learning in managers' choice of formation,'' which names lineage transmission along mentor-protege ties as open work, I paused before running those confirmatory tests to register them here. That paper was itself preregistered (osf.io/er4dx). The account that follows states what the preliminary work had established, so hindsight cannot blur the confirmatory line.

This preregistration covers new analyses of a dataset that already exists and that we have already explored in a preliminary, proof-of-concept phase. Integrity requires stating what that preliminary work established, so the confirmatory predictions below are understood as motivated by it but not yet formally tested.

The preliminary phase built the four coaching genes (offensive aggression, shotgun, tempo, defensive aggression) and their sub-components for 2006-2024, offensive aggression, tempo, and defensive aggression being composites while shotgun is a single measure, with the defensive-aggression genes from 2016 and man coverage from 2018, and examined as a proof of concept how they transmit from coordinator to head coach and how they relate to Coach WAR (wins above replacement). That preliminary read used ordinary frequentist correlations, and it suggested a qualitative pattern: the schematic-identity traits (shotgun, tempo) appeared to transmit from coordinator to head coach more strongly than they predicted winning, whereas offensive aggression, an approach-to-scenarios trait, appeared to predict Coach WAR more strongly than it transmitted. Those impressions motivated the hypotheses below. The genes have since been rebuilt with the frozen, leakage-free pipeline, in which each play's expectation is scored by a model that never trained on that coach and the cross-fitting folds are balanced across eras. We have not computed or inspected any transmission ($h^{2}$) or selection (S) quantity on these rebuilt genes; the impressions above rest only on the earlier frequentist correlations, so no preliminary point estimate is carried forward and the confirmatory line stays clean.

Three analyses are novel and preregistered here. C1 estimates transmission ($h^{2}$) as the parent-offspring coefficient of a Bayesian multilevel regression, in place of a raw correlation. C2 estimates per-trait repeatability as a variance-components ICC, tabulated as the $h^{2}$ ceiling. C3 extends $h^{2}$ and S from the four genes to the ten sub-traits and tests the cross-trait $h^{2}$-versus-S relationship. None of the confirmatory quantities has been computed on the frozen pipeline.

We commit to not inspecting any C1-C3 output until this document is posted.

\section*{1. Confirmatory hypotheses}

Each trait is a standardized (z-scored) coach-season phenotype: the deviation of observed on-field behavior from a play-level expectation model, attributed to the head coach who led those games. The offensive, shotgun, and tempo genes are measured on the plays of the coach's own offense, and the defensive-aggression genes on the team's defensive plays; both are attributed to the head coach. Every play is assigned to the coach who led the team for that game, so a within-season head-coaching change splits the season between the two coaches at the game of the change. Coordinators are not assigned their own gene; the phenotype is defined only at the head-coach level, so a play is attributed to one coach and never to both a head coach and their coordinator at once. A coordinator-to-head-coach transition therefore pairs two different coaches: the outcome is the protege's own phenotype once they became a head coach, and the predictor is the phenotype of the head coach they apprenticed under, measured on that mentor's team over the seasons the protege served there as a coordinator. All inference is contemporary-group (within-season) adjusted so that shared league-wide drift is not counted as transmission or selection (Section 3). $h^{2}$ is transmissibility, the mentor-to-protege resemblance in a trait (Section 3), and S is the selection index, a trait's correlation with Coach WAR.

H1 (approach is selected, identity is inherited). Coaching traits split into two families that behave differently. The approach-to-scenarios traits, offensive aggression (going for it on fourth down, pass-heavy play-calling, beyond-the-sticks downfield passing, two-point attempts) and defensive aggression (loading the box, sending extra rushers, playing man vs. zone), are more strongly selected, meaning they predict Coach WAR, than the schematic-identity traits. The schematic-identity traits, shotgun and tempo, are more strongly inherited, meaning they transmit down the tree, than the approach traits, and their apparent link to winning disappears under contemporary-group adjustment.

H2 (repeatability is not heritability). Offensive aggression is a repeatable individual trait that is nonetheless not heritable; its within-coach repeatability clearly exceeds its heritability.

H3 (cross-trait relationship). Across the ten sub-traits, heritability and selection are negatively related; the traits that transmit most strongly are not the ones that predict winning.

\section*{2. Frozen trait set}

The primary cross-trait test uses ten sub-traits, frozen here.

\begin{center}
\begin{tabular}{llll}
\toprule
Gene & Sub-trait & Transition tested & Data window \\
\midrule
Offensive aggression & fourth\_down & OC-\textgreater{}HC & 2006-2024 \\
Offensive aggression & pass\_heavy & OC-\textgreater{}HC & 2006-2024 \\
Offensive aggression & deep\_pass & OC-\textgreater{}HC & 2006-2024 \\
Offensive aggression & two\_point & OC-\textgreater{}HC & 2006-2024 \\
Tempo & no\_huddle & OC-\textgreater{}HC & 2006-2024 \\
Tempo & pace & OC-\textgreater{}HC & 2006-2024 \\
Defensive aggression & box\_stacking & DC-\textgreater{}HC & 2016-2024 \\
Defensive aggression & pass\_rush & DC-\textgreater{}HC & 2016-2024 \\
Defensive aggression & man\_coverage & DC-\textgreater{}HC & 2018-2024 \\
Shotgun & shotgun & OC-\textgreater{}HC & 2006-2024 \\
\bottomrule
\end{tabular}
\end{center}

Three composites, offensive aggression, tempo, and defensive aggression, appear on the $h^{2}$-versus-S map as labeled exemplars but are not counted as independent points in the cross-trait test; shotgun is already a single sub-trait.

All ten sub-traits enter the primary cross-trait test.

Each coach contributes at most one transition per coordinator role: the most recent coordinator stint that precedes a head-coaching stint, paired with the first head-coaching stint beginning after it. Coaches with more than one coordinator-then-head-coach arc are represented by their most recent arc, which avoids pseudo-replication.

\section*{3. Analysis plan}

Analyses are in Python, with the Bayesian models in PyMC to keep the whole pipeline in one language and a single-command reproducible bundle. A frequentist ordinary-least-squares fit of the same regression, with mentor-clustered standard errors, is reported only as a sanity cross-check.

Every trait is contemporary-group adjusted (within-season demeaning) before modeling. Because each phenotype is expressed as a deviation from its own season's league mean, shared league-wide drift is removed exactly once; this within-season adjustment is the single era control, and no further season or era term enters the models, so transmission and repeatability are net of era with no double correction. Coach WAR is already a within-season, replacement-relative phenotype, with a between-season variance share of about 3\%, and receives no further era adjustment.

Each gene is an estimate with a known per-coach-season sampling variance, from a reliability decomposition applied to every sub-trait. Because measurement noise attenuates both the regression slope, through error in the mentor-side predictor, and the repeatability variance ratio, and does so unevenly across traits of differing reliability, C1 and C2 carry this sampling variance as an explicit measurement-error level, so that $h^{2}$ and repeatability are estimated for the latent trait rather than its noisy observation. Where a per-observation variance is not cleanly estimable for a sub-trait, the estimate is disattenuated by its reliability and both raw and disattenuated values are reported. Reliability is estimated and reported for every sub-trait. Selection S is inverse-variance weighted, so the cross-trait test compares $h^{2}$ and S on the same disattenuated footing.

For each trait, C1 estimates transmissibility ($h^{2}$) as the coefficient of a Bayesian multilevel parent-offspring regression, the estimator matched to this design: both the mentor's and the protege's phenotypes are measured directly, so transmission is read from their resemblance rather than inferred from a relatedness matrix that could only be assumed. Concretely, the protege's own head-coaching trait is regressed on the mentor's trait, the latter measured on the mentor's team over the seasons the protege served there as a coordinator, and the standardized slope is $h^{2}$, reported as a posterior median with a 95\% HDI. Because a mentor transmits a whole trait rather than half a genome, the slope itself is the transmissibility, with no Mendelian doubling. Because the mentor's and protege's phenotypes are each already expressed relative to their own season, the slope is net of shared league drift with no additional era term, and the mentor-side value enters as a latent variable carrying its known measurement variance, an errors-in-variables correction for the downward attenuation that predictor noise would otherwise cause. Priors are Normal(0, 1) on the fixed effects (slope and intercept; traits are z-scored) and half-Normal(0, 1) on the residual standard deviation, with a half-Normal(0, 1) prior on the mentor standard deviation in the varying-intercept robustness fit. Because era enters only through the within-season demeaning of Section 3, no era standard deviation is estimated. Sampling uses 4 chains with 2000 post-warmup draws each; convergence targets R-hat at or below 1.01 and bulk effective sample size above 400, using a non-centered parameterization. Proteges who trained under the same mentor are not independent, so as a robustness check the slope is refit with a mentor varying intercept and, in a frequentist version, with mentor-clustered standard errors; the estimate should be stable. No pedigree animal model is fit (Section 6).

Selection S is the era-adjusted, inverse-variance-weighted correlation of a trait with Coach WAR. Sub-trait S is computed with the identical estimator already used for the composites and the defensive-aggression components, a mechanical extension pre-specified here.

C2 estimates a variance-components ICC (coach random effect) per trait, net of era, as a posterior median with a 95\% HDI, tabulated against $h^{2}$. Because heritability cannot exceed repeatability, the two are compared on their posterior draws: for each trait we report P($h^{2}$ \textgreater{} repeatability), and a value not near zero flags a specification problem to resolve before interpretation rather than a substantive result.

C3 assembles the ($h^{2}$, S) pair for the ten sub-traits and computes the Spearman rank correlation between them with a 95\% bootstrap interval. This cross-trait rank correlation between a trait's heritability and its association with fitness is the established test of the heritability-fitness relationship in quantitative genetics (Mousseau and Roff 1987; Kruuk et al. 2000, who report Spearman correlations near -0.76 across a comparably small set of traits); the rank form is used because heritability estimates are noisy and non-normally distributed, and we apply it here to a cultural transmission system. The ten sub-traits index distinct on-field behaviors, for example fourth-down decisions versus deep targeting, rather than re-scalings of one measure, and the full inter-sub-trait correlation matrix is reported alongside the result so readers can judge their independence.

The decision rules are one per hypothesis. H1 is supported if, comparing the approach family (offensive aggression, defensive aggression) with the identity family (shotgun, tempo), the approach family's mean S exceeds the identity family's and the identity family's mean $h^{2}$ exceeds the approach family's, each difference with a 95\% interval excluding 0. H2 is supported if the quantity (aggression repeatability minus aggression $h^{2}$) has a 95\% HDI excluding 0. H3 is supported if the Spearman rank correlation between $h^{2}$ and S across the sub-traits is negative with a 95\% bootstrap interval excluding 0.

Each hypothesis is judged by its single pre-specified rule above; we report the full posterior or bootstrap interval alongside each outcome so that magnitude and uncertainty are visible, not only the accept-or-reject decision.

\section*{4. Confirmatory vs exploratory boundary}

The confirmatory core is C1 (per-trait $h^{2}$), C2 (repeatability versus $h^{2}$), C3 (the cross-trait $h^{2}$-versus-S relationship), and the H1 approach-versus-identity family contrast. The genes feeding these tests come from the leakage-free, coach-held-out pipeline. Any analysis, contrast, covariate, or decision rule not specified in this document is exploratory, not confirmatory; only the tests and rules stated here are confirmatory.

The following analyses are declared exploratory and are corrected among themselves under a Benjamini-Hochberg FDR, separate from the confirmatory tests, which are judged by their pre-specified interval rules. The first is an era-residualized Moran's I on the mentor network, using an undirected, row-normalized adjacency of mentor-protege ties, with a within-era permutation null that takes node era as career midpoint and uses at least 1000 within-era shuffles. The second is reproductive fitness as an exposure-normalized rate, offspring (proteges who themselves became head coaches) per coordinator-season with a minimum-exposure floor, together with the test of whether coach quality predicts that rate net of exposure.

\section*{5. Data and code availability}

Everything is public, seed-fixed, and runnable end-to-end: this repository plus the sibling Coach\_WAR repository, which is the source of the WAR phenotype. Selection S is computed on Coach WAR at commit dffe2f1 (2026-06-18). If the WAR estimator is revised, for example in response to methodological review, S is recomputed on the revised version as a pre-specified robustness check and both are reported; the confirmatory predictions concern the sign of the contrasts, not the exact value of S.

\section*{6. Limitations}

The transmission samples are small, about 40 for OC-\textgreater{}HC and about 17 for DC-\textgreater{}HC, so the heritability posteriors will be wide; the main result rests on the cross-trait pattern (H3) and on the robustness of each slope across specifications, not on any single precise $h^{2}$. Coaching pedigrees are also non-tree, with multiple simultaneous mentors and overlapping careers, and coaching relatedness can only be assumed by analogy rather than derived from a Mendelian process; we therefore do not fit a pedigree relatedness model and rely on the parent-offspring regression, with a mentor varying intercept only as a non-independence check. Genes attribute team play-calling to the head coach and do not isolate a coordinator's individual contribution within a staff, so the transmission tests do not measure a coordinator's own phenotype. They instead ask whether the head-coaching style a coach apprenticed under, the phenotype of the head coach and team over their overlapping seasons, predicts that coach's own style once they became a head coach. Such mentor-protege resemblance is consistent with social learning but cannot on its own be separated from assortative hiring, in which head coaches recruit coordinators who already share their approach; we therefore read $h^{2}$ as vertical transmissibility in the broad sense rather than as social learning specifically.

\end{document}
